% Syntra LaTeX Template
% For Arabic support, use XeLaTeX with:
%   \usepackage{fontspec}
%   \usepackage{polyglossia}
%   \setdefaultlanguage{english}
%   \setotherlanguage{arabic}
%   \newfontfamily\arabicfont[Script=Arabic]{Amiri}

\documentclass[12pt,a4paper]{article}

% Font and encoding
\usepackage[utf8]{inputenc}
\usepackage[T1]{fontenc}

% Math packages
\usepackage{amsmath,amssymb}

% Table support
\usepackage{longtable,booktabs}

% Graphics and links
\usepackage{graphicx}
\usepackage{hyperref}
\usepackage{geometry}

% Page geometry
\geometry{
    left=2.5cm,
    right=2.5cm,
    top=3cm,
    bottom=3cm
}

% Fix pandoc tightlist issue
\providecommand{\tightlist}{\setlength{\itemsep}{0pt}\setlength{\parskip}{0pt}}

% Hyperref settings
\hypersetup{
    colorlinks=true,
    linkcolor=blue,
    filecolor=magenta,
    urlcolor=cyan,
    pdftitle={Syntra Research Documentation},
    pdfauthor={Syntra-Env},
    pdfpagemode=FullScreen
}

\begin{document}

\section{References}\label{references}

Academic and research sources that inform and inspire what
we\textquotesingle re building.

\subsection{Value Alignment and
Coherence}\label{value-alignment-and-coherence}

\begin{itemize}
\tightlist
\item
  {[}ref-1{]} Anthropic, "Claude\textquotesingle s Constitution,"
  \emph{Anthropic}, Jan. 2024. {[}Online{]}. Available:
  https://www.anthropic.com/news/claude-new-constitution
\item
  {[}ref-2{]} A. Bai, et al., "Constitutional AI: Harmlessness from AI
  feedback," \emph{Anthropic}, 2022. {[}Online{]}. Available:
  https://www.anthropic.com/index.php/constitutional-ai
\end{itemize}

\subsection{Architecture and
Philosophy}\label{architecture-and-philosophy}

\begin{itemize}
\tightlist
\item
  {[}ref-3{]} S. A. Senchal, "The God Conjecture: Mapping theological
  creation narratives to computational physics via the Ruliad
  framework," 2024. {[}Online{]}. Available:
  https://github.com/SASenchal/God-Conjecture
\end{itemize}

\subsection{Geometric and Mathematical
Foundations}\label{geometric-and-mathematical-foundations}

\begin{itemize}
\tightlist
\item
  {[}ref-4{]} N. Spivack, "Toward a geometric theory of information
  processing: A research program," \emph{Nova Spivack}, May 2025.
  {[}Online{]}. Available:
  https://www.novaspivack.com/science/toward-a-geometric-theory-of-information-processing-a-research-program
\item
  {[}ref-5{]} N. Spivack, "Information processing complexity as
  spacetime curvature: A formal derivation and physical unification,"
  \emph{Nova Spivack}, June 2025. {[}Online{]}. Available:
  https://www.novaspivack.com/science/information-processing-complexity-as-spacetime-curvature-a-formal-derivation-and-physical-unification
\item
  {[}ref-6{]} J. Bürge, "Neural geodesic flows: Time series prediction
  with a learnable geodesic flow," \emph{GitHub}, 2025. {[}Online{]}.
  Available: https://github.com/julianbuerge/neural\_geodesic\_flows
\end{itemize}

\subsection{Holonomic Unified Field Dynamics
(HUFD)}\label{holonomic-unified-field-dynamics-hufd}

\begin{itemize}
\tightlist
\item
  {[}ref-7{]} J. Harlow, "Holonomic unified field dynamics (HUFD): A
  constant-attention manifold for recurrent field computation,"
  \emph{Local Reference}. Available:
  \texttt{docs/Information\ Geometry/jaim\ harlow/}
\item
  {[}ref-8{]} J. Harlow, "Retrieval and reasoning in HUFD: From constant
  attention to holonomic geometry," \emph{Local Reference}. Available:
  \texttt{docs/Information\ Geometry/jaim\ harlow/}
\end{itemize}

\subsection{Categorical and Information-Theoretic
Foundations}\label{categorical-and-information-theoretic-foundations}

\begin{itemize}
\tightlist
\item
  {[}ref-9{]} UOR Foundation, "UOR Framework: Universal Object
  Reference," \emph{GitHub}, 2026. {[}Online{]}. Available:
  https://github.com/UOR-Foundation/UOR-Framework
\item
  {[}ref-10{]} I. Paveliev, "Universal Object Reference (UOR) Coordinate
  System," \emph{LinkedIn}, 2025. {[}Online{]}. Available:
  https://www.linkedin.com/posts/trinityinvestor\_introducing-the-universal-object-reference-activity-7417508099625938945-71d8
\end{itemize}

\subsection{Researchers and
Practitioners}\label{researchers-and-practitioners}

\begin{itemize}
\tightlist
\item
  {[}ref-11{]} J. Harlow, "Non-Abelian holonomy, quantum computing,
  high-energy physics, photonics, and HUFD architecture," \emph{Google
  Scholar}. {[}Online{]}. Available:
  https://scholar.google.com/citations?user=\_6B\_TlcAAAAJ\&hl=en
\end{itemize}

\subsection{Core Concepts (Glossary)}\label{core-concepts-glossary}

{\def\LTcaptype{none} % do not increment counter
\begin{longtable}[]{@{}ll@{}}
\toprule\noalign{}
Term & Definition \\
\midrule\noalign{}
\endhead
\bottomrule\noalign{}
\endlastfoot
\textbf{Holonomy} & The rotation of an internal state (meaning) as it is
transported along a path in the text. \\
\textbf{Gauge Stability} & The mathematical consistency of features
across different textual contexts. \\
\textbf{Prime Framework} & Representing coordinates as prime
factorizations to ensure observer invariance. \\
\textbf{Semantic Tension} & Local curvature in the manifold representing
high interpretive demand or anomaly. \\
\textbf{NCU} & Nafs-Centric Universe; a framework for tracing
perspective and agency in narratives. \\
\end{longtable}
}

\end{document}
